% Part: computability
% Chapter: computability-theory
% Section: prop-reduce

\documentclass[../../../include/open-logic-section]{subfiles}

\begin{document}

\olfileid{cmp}{thy}{ppr}
\olsection{خواص قابلية الاختزال}

نكتب $A\leq_m B$ إذا كانت $A$ تختزل إلى~$B$. ويوحي هذا الترميز بأنه
إذا كان $A\leq_m B$، فإن $A$ «ليست أصعب من»~$B$، وإن $B$ «لا تقل
صعوبة عن»~$A$، وإن $\le_m$ ترتب المجموعات، أو مسائل القرار، بحسب
تعقيدها مثلما ترتب $\le$ المعتادة الأعداد. وتؤيد القضيتان الآتيتان هذا
الحدس. تقول الأولى إن $\le_m$ متعدية.

\begin{prop}
\ollabel{prop:trans-red}
إذا كان $A\leq_m B$ و$B\leq_m C$، فإن $A\leq_m C$.
\end{prop}

\begin{proof}
ينتج من تركيب اختزال $A$ إلى~$B$ مع اختزال $B$ إلى~$C$ اختزال $A$
إلى~$C$.
\end{proof}

\begin{prob}
أثبت \olref{prop:trans-red} ببيان أنه إذا كانت $f$ و~$g$ اختزالين
متعددين إلى واحد لـ$A$ إلى~$B$ ولـ$B$ إلى~$C$، على الترتيب، فإن
$\comp{f}{g}$ اختزال متعدد إلى واحد لـ$A$ إلى~$C$.
\end{prob}

\begin{prop}
\ollabel{prop:reduce}
لتكن $A$ و$B$ أي مجموعتين، وافترض أن $A\leq_m B$.
\begin{enumerate}
\item إذا كانت $B$ قابلة للتعداد حسابيًا، كانت~$A$ كذلك.
\item إذا كانت $B$ قابلة للحساب، كانت~$A$ كذلك.
\end{enumerate}
\end{prop}

\begin{proof}
لتكن $f$ اختزالًا متعددًا إلى واحد من $A$ إلى~$B$. ولإثبات الادعاء
الأول، يكفي التحقق من أنه إذا كانت $B$ مجال دالة جزئية~$g$، فإن $A$
مجال~$\comp{f}{g}$:
\begin{align*}
x \in A & \text{ إذا وفقط إذا كان } f(x) \in B \\
& \text{ إذا وفقط إذا كانت }  g(f(x)) \fdefined.
\end{align*}

ولإثبات الادعاء الثاني، تذكر أنه إذا كانت $B$ قابلة للحساب، كانت $B$
و~$\Complement{B}$ قابلتين للتعداد حسابيًا
(\olref[cmp]{thm:ce-comp}). وليس من الصعب التحقق من أن $f$ أيضًا
اختزال متعدد إلى واحد لـ$\Complement{A}$ إلى $\Complement{B}$؛ ولذلك
تكون $A$ و~$\Complement{A}$ !!{computably enumerable} بحسب الجزء الأول
من هذا البرهان. ومن ثم تكون $A$ قابلة للحساب أيضًا بحسب
\olref[cmp]{thm:ce-comp}. (وبديلًا من ذلك، يمكنك التحقق من أن
$\Char{A}=\comp{f}{\Char{B}}$؛ فإذا كانت $\Char{B}$ قابلة للحساب،
كانت~$\Char{A}$ كذلك.)
\end{proof}

\begin{prob}
افترض أن $f$ اختزال متعدد إلى واحد لـ$A$ إلى~$B$. بين أن $f$ أيضًا
اختزال متعدد إلى واحد لـ$\Complement{A}$ إلى $\Complement{B}$.
\end{prob}

\begin{prob}
بين أنه إذا كانت $f$ اختزالًا متعددًا إلى واحد لـ$A$ إلى~$B$، فإن
$\Char{A}=\comp{f}{\Char{B}}$.
\end{prob}

\begin{digress}
يفيد في سياقات أخرى، ولا سيما في إثبات نتائج عدم القابلية للقرار، مفهوم
أعم لقابلية الاختزال يسمى \emph{قابلية الاختزال بتورنغ}. لاحظ أنه بحسب
\olref[cmp]{cor:comp-k} لا تقبل متممة~$K_0$ الاختزال إلى~$K_0$، لأنها
غير قابلة للتعداد حسابيًا. لكنك، بحسب الحدس، لو عرفت أجوبة الأسئلة عن
$K_0$، لعرفت أجوبة الأسئلة عن متممتها أيضًا. ويقال إن مجموعة $A$
قابلة للاختزال بتورنغ إلى~$B$ إذا أمكن تعيين أجوبة الأسئلة عن~$A$
بإجراء قابل للحساب يستطيع طرح أسئلة عن~$B$. وهذا أوسع من قابلية
الاختزال متعددًا إلى واحد؛ إذ لا يسمح في الأخيرة إلا (1) بطرح سؤال
واحد عن $B$، و(2) يجب أن تترجم إجابة «نعم» إلى إجابة «نعم» عن السؤال
المتعلق بـ$A$، وكذلك «لا». ويظل صحيحًا أنه إذا كانت $A$ قابلة للاختزال
بتورنغ إلى~$B$ وكانت $B$ قابلة للحساب، فإن $A$ قابلة للحساب أيضًا
(مع أن العبارة المناظرة لا تصح، كما رأينا، في القابلية للتعداد حسابيًا).

ينبغي أن تفكر في مختلف مفاهيم قابلية الاختزال التي ناقشناها، وأن تفهم
الفروق بينها. غير أننا لن نتناول في هذا الفصل إلا قابلية الاختزال
متعددًا إلى واحد. ومن عَرَض القول أن لنوعي الاختزال المناقشين في الفقرة
السابقة نظيرين في التعقيد الحسابي، مع الشرط الإضافي القائل بأن تعمل
آلات تورنغ في زمن كثير الحدود: فنسخة التعقيد من قابلية الاختزال متعددًا
إلى واحد تعرف بـ\emph{قابلية الاختزال بكارب}، أما نسخة التعقيد من
قابلية الاختزال بتورنغ فتعرف بـ\emph{قابلية الاختزال بكوك}.
\end{digress}

\end{document}
